\subsubsection*{NAND}

% BEGIN GENERATED ARTIFACT: def:nand-connective-propositional-logic
\begin{definitionbox}{Definition (NAND Connective)}
\begin{definition}[NAND Connective]
\label{def:nand-connective-propositional-logic}
The \emph{NAND} connective, written \(\uparrow\), is the binary connective
defined by
\[
\varphi\uparrow\psi\equiv\neg(\varphi\land\psi).
\]
\end{definition}
\end{definitionbox}

\begin{remark*}[Standard quantified statement]
\[
\widehat v(\varphi\uparrow\psi)=\True
\Longleftrightarrow
\text{not both }\widehat v(\varphi)=\True\text{ and }\widehat v(\psi)=\True.
\]
\end{remark*}

\begin{remark*}[Predicate reading]
\(\operatorname{NAND}(\varphi,\psi)\) means the negation of the conjunction of
\(\varphi\) and \(\psi\).
\end{remark*}

\begin{remark*}[Interpretation]
NAND is read as "not both." It is false only when both inputs are true. This
single connective can reproduce negation, conjunction, and disjunction.
\end{remark*}

\begin{dependencies}
\begin{itemize}
  \item \hyperref[def:logical-connective]{Logical Connective}
  \item \hyperref[def:truth-assignment-propositional-logic]{Truth Assignment}
\end{itemize}
\end{dependencies}
% END GENERATED ARTIFACT: def:nand-connective-propositional-logic

% BEGIN GENERATED ARTIFACT: def:nand-only-formula-propositional-logic
\begin{definitionbox}{Definition (NAND-Only Formula)}
\begin{definition}[NAND-Only Formula]
\label{def:nand-only-formula-propositional-logic}
A \emph{NAND-only formula} is a well-formed formula whose only connective is
\(\uparrow\).
\end{definition}
\end{definitionbox}

\begin{remark*}[Standard quantified statement]
\[
\varphi\text{ is NAND-only}
\Longleftrightarrow
\varphi\in\WFF_{\{\uparrow\}}.
\]
\end{remark*}

\begin{remark*}[Predicate reading]
\(\operatorname{NANDOnly}(\varphi)\) means that every connective occurrence in
\(\varphi\) is an occurrence of \(\uparrow\).
\end{remark*}

\begin{remark*}[Interpretation]
A NAND-only formula is syntactically restricted but expressively powerful. The
restriction is severe at the level of primitive symbols, but not at the level of
truth functions.
\end{remark*}

\begin{dependencies}
\begin{itemize}
  \item \hyperref[def:nand-connective-propositional-logic]{NAND Connective}
  \item \hyperref[def:well-formed-formula]{Well-Formed Formula}
\end{itemize}
\end{dependencies}
% END GENERATED ARTIFACT: def:nand-only-formula-propositional-logic

% BEGIN GENERATED ARTIFACT: thm:nand-functionally-complete-propositional-logic
\begin{theorembox}{Theorem (NAND is Functionally Complete)}
\begin{theorem}[NAND is Functionally Complete]
\label{thm:nand-functionally-complete-propositional-logic}
The one-connective basis \(\{\uparrow\}\) is functionally complete.

\noindent\hyperref[prf:nand-functionally-complete-propositional-logic]{Go to proof.}
\end{theorem}
\end{theorembox}

\begin{remark*}[Standard quantified statement]
\[
\operatorname{FunctionallyComplete}(\{\uparrow\}).
\]
\end{remark*}

\begin{remark*}[Predicate reading]
Every Boolean function can be represented by a NAND-only formula.
\end{remark*}

\begin{remark*}[Interpretation]
NAND can define negation and conjunction:
\[
\neg\varphi\equiv\varphi\uparrow\varphi,
\qquad
\varphi\land\psi\equiv(\varphi\uparrow\psi)\uparrow(\varphi\uparrow\psi).
\]
Once negation and conjunction are definable, functional completeness follows
from the functional completeness of \(\{\neg,\land\}\).
\end{remark*}

\begin{dependencies}
\begin{itemize}
  \item \hyperref[def:nand-connective-propositional-logic]{NAND Connective}
  \item \hyperref[def:nand-only-formula-propositional-logic]{NAND-Only Formula}
  \item \hyperref[def:functionally-complete-connective-set-propositional-logic]{Functionally Complete Connective Set}
  \item \hyperref[thm:negation-conjunction-functionally-complete-propositional-logic]{Negation and Conjunction are Functionally Complete}
\end{itemize}
\end{dependencies}
% END GENERATED ARTIFACT: thm:nand-functionally-complete-propositional-logic
